\chapter{Diferenciación, espectro y cálculo holomorfo}\label{cap:v3-11}

La inversión de matrices es una operación no lineal. Sin embargo, lejos de las matrices singulares, una perturbación pequeña admite una expansión controlada. El mismo mecanismo funciona en cualquier álgebra de Banach: la serie geométrica proporciona primero una derivada, después una función resolvente y finalmente una manera de evaluar funciones holomorfas en elementos del álgebra. Desarrollaremos esta cadena sin confundir el producto no conmutativo con el producto de números.

\section{La derivada como aproximación lineal}\label{lec:v3-21}
\fuentelectura{21}{La derivada de Fréchet}

\begin{definicion}
Sean $X,Y$ espacios normados sobre $\F$, $U\subset X$ abierto y $a\in U$. Una aplicación $f:U\to Y$ es \emph{diferenciable en el sentido de Fréchet} en $a$ si existe $T\in\Lcal(X,Y)$ tal que
\[
 f(a+h)=f(a)+Th+r(h),\qquad
 \frac{\norm{r(h)}}{\norm h}\longrightarrow0\quad(h\to0,\ h\ne0).
\]
Escribimos $Df(a)=T$. Cuando $\F=\C$, la aplicación $T$ debe ser compleja lineal; la diferenciabilidad real es una condición diferente.
\end{definicion}
\index{derivada de Fréchet}

El cociente controla el error simultáneamente en todas las direcciones. Las derivadas direccionales aisladas no proporcionan este control uniforme. La apertura de $U$ garantiza que todos los incrementos suficientemente pequeños están permitidos.

\begin{proposicion}[Reglas elementales]\label{v3:reglas-frechet}
La derivada, cuando existe, es única. La diferenciabilidad implica continuidad; la suma y la multiplicación por escalares se diferencian término a término. Si $f:U\to Y$ es diferenciable en $a$ y $g:V\to Z$ lo es en $f(a)$, con $f(U)\subset V$, entonces
\[
 D(g\circ f)(a)=Dg(f(a))\,Df(a).
\]
\end{proposicion}
\begin{proof}
Si $T,S$ son derivadas, sustituyendo $h=tv$ y dividiendo entre $|t|$ obtenemos $\norm{(T-S)v}=0$ para cada $v$. Si $r(h)=o(\norm h)$, para $h$ pequeño se tiene $\norm{f(a+h)-f(a)}\le(\norm T+1)\norm h$, lo cual prueba continuidad incluso cuando $T=0$. Las reglas de linealidad se siguen de que una combinación lineal de errores $o(\norm h)$ vuelve a ser de ese orden.

Para la composición pongamos $k=f(a+h)-f(a)=Th+r(h)$. Entonces $\norm k\le C\norm h$ y
\[
 g(f(a)+k)-g(f(a))=Sk+s(k)=STh+Sr(h)+s(k).
\]
El primer resto es $o(\norm h)$. Para el segundo, si $k\ne0$,
$\norm{s(k)}/\norm h=(\norm{s(k)}/\norm k)(\norm k/\norm h)\to0$; si $k=0$, el resto es cero. Esta observación evita dividir entre un incremento que puede anularse.
\end{proof}

\begin{ejemplo}
Una aplicación lineal acotada $T$ tiene derivada $T$ en todos los puntos. En un espacio de Hilbert real, para $b(x,y)=\ip{x}{y}$,
\[
 Db(x,y)[h,k]=\ip h y+\ip x k,
\]
pues el resto $\ip h k$ tiene módulo a lo sumo $\norm h\norm k$. Usamos en el producto la norma $\norm{(h,k)}=\norm h+\norm k$. El caso complejo exige distinguir las dos nociones de derivada; véase el ejercicio~\ref{ej:v3-l21-2}.
\end{ejemplo}

\begin{proposicion}
Si $f:X\to Y$ es diferenciable y $Df(x)=0$ para todo $x$, entonces $f$ es constante. Más generalmente, basta que el dominio sea un abierto conexo.
\end{proposicion}
\begin{proof}
Fijemos un segmento $[a,b]$ contenido en el dominio y $\varphi\in Y^*$. La función $t\mapsto\varphi(f(a+t(b-a)))$ tiene derivada cero en $0<t<1$. El teorema del valor medio, aplicado a sus partes real e imaginaria si es necesario, da $\varphi(f(b)-f(a))=0$. Hahn--Banach separa puntos y por ello $f(a)=f(b)$. En un abierto conexo, el conjunto de puntos que pueden unirse a $a$ por una poligonal contenida en el abierto es abierto; su complemento también lo es, pues cada punto tiene una bola convexa contenida en el dominio. La conexidad implica que ese conjunto es todo el dominio. Aplicamos el argumento en cada segmento de la poligonal.
\end{proof}

\subsection{Producto, potencias e inversión}
En esta sección $A$ es un álgebra de Banach unitaria no nula. Podemos normalizar $\norm1=1$ mediante la norma equivalente del ejercicio~\ref{ej:v3-l19-3}; esto no modifica diferenciabilidad, espectro ni límites de raíces de normas de potencias.

\begin{proposicion}\label{v3:derivada-producto}
La multiplicación $m:A\times A\to A$ es diferenciable y
\[
 Dm(x,y)[h,k]=hy+xk.
\]
\end{proposicion}
\begin{proof}
Al desarrollar $(x+h)(y+k)$ obtenemos $xy+hy+xk+hk$. La parte propuesta es lineal y acotada, y $\norm{hk}\le\norm h\norm k\le\norm{(h,k)}^2/4$. El orden $hy$ es esencial: en un álgebra no conmutativa no puede sustituirse por $yh$.
\end{proof}

\begin{lema}[Límite submultiplicativo]\label{v3:limite-raices}
Para cada $x\in A$ existe
\[
 \rho_0(x):=\lim_{n\to\infty}\norm{x^n}^{1/n}
 =\inf_{n\ge1}\norm{x^n}^{1/n}\le\norm x.
\]
Si $\rho_0(x)<1$, la serie $\sum_{n\ge0}x^n$ converge en norma y es la inversa de $1-x$.
\end{lema}
\begin{proof}
Si alguna potencia es cero, todas las posteriores lo son y el resultado es inmediato. En otro caso fijemos $N$. Para $n=qN+r$, $0\le r<N$, se tiene
$\norm{x^n}\le\norm{x^N}^{q}C_N$, donde $C_N=\max_{0\le r<N}\norm{x^r}$. Por tanto
$\limsup_n\norm{x^n}^{1/n}\le\norm{x^N}^{1/N}$. Tomando el ínfimo en $N$ y comparando con el límite inferior queda probada la igualdad. Si $\rho_0(x)<s<1$, para $n$ grande $\norm{x^n}\le s^n$, de modo que la serie converge absolutamente. Multiplicar sus sumas parciales por $1-x$ a derecha e izquierda da $1-x^{n+1}$; la continuidad del producto permite pasar al límite.
\end{proof}

\begin{teorema}[Apertura del grupo de invertibles]\label{v3:derivada-inversa}
El conjunto $G(A)$ es un grupo multiplicativo abierto. Si $x\in G(A)$ y $\norm h<\norm{x^{-1}}^{-1}$, entonces $x+h$ es invertible. La aplicación $\operatorname{inv}:G(A)\to G(A)$ es diferenciable y
\[
 D\operatorname{inv}(x)[h]=-x^{-1}hx^{-1}.
\]
\end{teorema}
\begin{proof}
Las identidades $(xy)^{-1}=y^{-1}x^{-1}$ y $(x^{-1})^{-1}=x$ prueban la afirmación algebraica. Como $x+h=x(1+x^{-1}h)$, la serie de Neumann da
\[
 (x+h)^{-1}=\sum_{n=0}^{\infty}(-x^{-1}h)^nx^{-1}
 =x^{-1}-x^{-1}hx^{-1}+R(h),
\]
con
\[
 \norm{R(h)}\le
 \frac{\norm{x^{-1}}^3\norm h^2}{1-\norm{x^{-1}}\norm h}.
\]
La cota prueba el resto requerido y la apertura. También demuestra continuidad de la inversión. En particular, la derivada depende continuamente de $x$, porque los factores $x^{-1}$ lo hacen.
\end{proof}

\section{Ejercicios sobre diferenciación y exponenciales}

\begin{ejercicio}\label{ej:v3-l21-1}
Demuestre la unicidad de la derivada de Fréchet. Explique por qué la hipótesis de que el dominio sea abierto es necesaria para esa afirmación.
\end{ejercicio}


\begin{ejercicio}\label{ej:v3-l21-2}
Determine si el producto interno de un espacio de Hilbert complejo, visto como aplicación $H\times H\to\C$, es diferenciable.
\end{ejercicio}


\begin{ejercicio}\label{ej:v3-l21-3}
Sean $x,y\in A$. Pruebe: (i) la serie $e^x=\sum_{n\ge0}x^n/n!$ converge absolutamente; (ii) si $xy=yx$, entonces $e^{x+y}=e^xe^y$; (iii) $e^x$ es invertible y su inversa es $e^{-x}$; (iv) $e^x=\lim_n(1+x/n)^n$; (v) la fórmula de producto de Trotter
\[
 e^{x+y}=\lim_{n\to\infty}(e^{x/n}e^{y/n})^n
\]
es válida sin suponer que $x$ e $y$ conmuten.
\end{ejercicio}


\begin{ejercicio}\label{ej:v3-l21-4}
Pruebe que $t\mapsto e^{tx}$ es diferenciable y su derivada es $xe^{tx}=e^{tx}x$.
\end{ejercicio}


\begin{ejercicio}\label{ej:v3-l21-5}
Sea $x\in\partial G(A)$. Pruebe: (i) para cada $y\in G(A)$,
$\norm{y^{-1}}\ge1/\norm{x-y}$; (ii) $x$ es un divisor topológico de cero por ambos lados, es decir, existen $u_n$ con $\norm{u_n}=1$ y $xu_n,u_nx\to0$; (iii) $x$ no tiene inversa izquierda ni derecha.
\end{ejercicio}


\section{El espectro en un álgebra de Banach}\label{lec:v3-22}
\fuentelectura{22}{El espectro de un elemento en un álgebra}

Desde ahora el campo es $\C$, salvo indicación contraria. Es indispensable esta hipótesis: una matriz real de rotación puede tener espectro real vacío.

\begin{definicion}
Para $x\in A$, su \emph{espectro} y su \emph{conjunto resolvente} son
\[
 \sigma_A(x)=\{\lambda\in\C:\lambda1-x\notin G(A)\},
 \qquad \varrho_A(x)=\C\setminus\sigma_A(x).
\]
Escribimos $R_x(\lambda)=(\lambda1-x)^{-1}$ para $\lambda\in\varrho_A(x)$.
\end{definicion}
\index{espectro!en un álgebra de Banach}

\begin{ejemplo}
Si $A=C(K)$, con $K$ compacto Hausdorff no vacío, entonces $\sigma(f)=f(K)$. En efecto, si $\lambda\notin f(K)$, la función $t\mapsto1/(\lambda-f(t))$ es continua y es la inversa de $\lambda-f$. Si $f(t_0)=\lambda$, evaluar cualquier supuesta identidad $(\lambda-f)g=1$ en $t_0$ da $0=1$.
\end{ejemplo}

\subsection{Integración vectorial y el puente con análisis complejo}
Para una función continua $F:[a,b]\to A$, las sumas de Riemann forman una red de Cauchy por continuidad uniforme; la completitud de $A$ define $\int_a^bF(t)\,dt$. La desigualdad triangular pasa al límite y da $\norm{\int F}\le\int\norm F$. Una curva a trozos $C^1$ define $\int_\gamma F(z)\,dz=\int F(\gamma(t))\gamma'(t)\,dt$. Toda aplicación lineal acotada conmuta con esta integral, porque conmuta con las sumas finitas y con sus límites. Las integrales dobles de funciones continuas sobre productos compactos pueden intercambiarse: las dos iteraciones son límites de las mismas sumas rectangulares.

Recordemos la forma precisa de los resultados escalares que utilizaremos. Si $f$ es holomorfa en un abierto $U$ y $\Gamma$ es un ciclo cerrado a trozos suave cuya región interior está contenida en $U$, entonces $\int_\Gamma f=0$. Para $w$ fuera de $\Gamma$,
\[
 \frac1{2\pi i}\int_\Gamma\frac{f(z)}{z-w}\,dz
 =\operatorname{Ind}(\Gamma,w)f(w)
\]
si $w\in U$; el miembro izquierdo es cero cuando el integrando es holomorfo en toda la región encerrada. En particular, para una curva simple orientada positivamente el índice interior vale uno. La simple holomorfía cerca de un punto, sin control de la región encerrada, no basta para el teorema integral.

Para recordar por qué no se necesita continuidad de la derivada en estas fórmulas, el argumento de Goursat subdivide un triángulo en cuatro y conserva uno cuya integral tiene al menos un cuarto del módulo original. Los triángulos anidados convergen a $z_0$; la expansión $f(z)=f(z_0)+f'(z_0)(z-z_0)+o(|z-z_0|)$ y la anulación de las integrales de las dos primeras partes dan, tras multiplicar por $4^n$, una cota que tiende a cero para la integral inicial. La cancelación de aristas interiores extiende la conclusión a regiones trianguladas. Al quitar un disco pequeño alrededor de $w$, se aplica el teorema integral a $f(z)/(z-w)$ en la región restante; en la circunferencia pequeña, separar $f(z)=f(w)+(f(z)-f(w))$ deja $2\pi i f(w)$ y un error acotado por $2\pi\sup|f(z)-f(w)|\to0$. La misma partición, contando orientaciones, da la versión para ciclos. Aproximar las fronteras a trozos suaves por polígonos dentro del abierto no cambia las integrales en el límite, por continuidad uniforme en una vecindad compacta de la curva.

Estas fórmulas valen también para funciones holomorfas con valores en $A$: aplíquense a $\varphi\circ F$ para cada $\varphi\in A^*$ y úsese Hahn--Banach para recuperar la identidad vectorial. Del mismo modo, una función entera vectorial acotada es constante. En el caso escalar esto se obtiene de la fórmula de Cauchy para la derivada, $|f'(w)|\le M/R$, dejando $R\to\infty$; la versión vectorial se obtiene por separación mediante funcionales.

\begin{teorema}[Compacidad y no vacuidad]\label{v3:espectro-algebra-compacto}
El espectro de cada elemento de un álgebra de Banach compleja unitaria no nula es compacto y no vacío; además $\sigma(x)\subset\{\lambda:|\lambda|\le\norm x\}$.
\end{teorema}
\begin{proof}
La apertura de $G(A)$ implica que el resolvente es abierto. Para $|\lambda|>\norm x$, la serie de Neumann da
\[
 R_x(\lambda)=\sum_{n=0}^{\infty}\frac{x^n}{\lambda^{n+1}},
 \qquad\norm{R_x(\lambda)}\le\frac1{|\lambda|-\norm x}.
\]
Por tanto el espectro es cerrado y acotado. La regla de la cadena y el teorema~\ref{v3:derivada-inversa} muestran que $R_x$ es holomorfa en el resolvente, con derivada $-R_x(\lambda)^2$. Si el espectro fuera vacío, esta función sería entera. La cota anterior la acota fuera de un disco y la continuidad la acota dentro. Liouville y el límite $R_x(\lambda)\to0$ en infinito implicarían $R_x\equiv0$, imposible porque $(\lambda1-x)R_x(\lambda)=1\ne0$.
\end{proof}

\begin{teorema}[Fórmula del radio espectral de Gelfand]\label{v3:radio-gelfand}
Para cada $x\in A$,
\[
 r(x):=\max_{\lambda\in\sigma(x)}|\lambda|
 =\lim_{n\to\infty}\norm{x^n}^{1/n}.
\]
\end{teorema}
\begin{proof}
Si $|\lambda|>\rho_0(x)$, el lema~\ref{v3:limite-raices} aplicado a $x/\lambda$ prueba que $\lambda1-x$ es invertible. Luego $r(x)\le\rho_0(x)$. Para la desigualdad opuesta fijemos $s>r(x)$. Sobre una circunferencia de radio $R>\norm x$, la expansión geométrica converge uniformemente y la integración término a término da
\[
 x^n=\frac1{2\pi i}\int_{|z|=R}z^nR_x(z)\,dz.
\]
El integrando es holomorfo en el anillo $s\le|z|\le R$; el teorema integral de Cauchy permite reemplazar $R$ por $s$. Por ello
\[
 \norm{x^n}\le s^{n+1}\max_{|z|=s}\norm{R_x(z)}.
\]
Tomar raíces $n$-ésimas y pasar al límite da $\rho_0(x)\le s$. Dejando $s\downarrow r(x)$, incluso cuando $r(x)=0$, obtenemos la igualdad.
\end{proof}

\begin{proposicion}[Operaciones elementales con el espectro]\label{v3:espectro-operaciones}
Para $x,y\in A$ y $\alpha\in\C$:
\[
 \sigma(x+\alpha1)=\sigma(x)+\alpha,\qquad
 \sigma(\alpha x)=\alpha\sigma(x).
\]
Si $x$ es invertible, $\sigma(x^{-1})=\{\lambda^{-1}:\lambda\in\sigma(x)\}$. Además,
\[
 \sigma(xy)\setminus\{0\}=\sigma(yx)\setminus\{0\}.
\]
\end{proposicion}
\begin{proof}
Las dos primeras fórmulas resultan de restar el parámetro espectral y factorizar el escalar; para $\alpha=0$, ambos lados son $\{0\}$ porque el espectro de $x$ es no vacío y $\sigma(0)=\{0\}$. Para $\mu\ne0$,
$x^{-1}-\mu1=-\mu x^{-1}(x-\mu^{-1}1)$; ambos factores externos son invertibles. El cero no está en ninguno de los espectros relevantes.

Si $\lambda\ne0$ y $xy-\lambda1$ es invertible, la identidad
\[
 (yx-\lambda1)^{-1}=\lambda^{-1}\bigl(y(xy-\lambda1)^{-1}x-1\bigr)
\]
se verifica multiplicando por ambos lados: use $(yx-\lambda1)y=y(xy-\lambda1)$ y $x(yx-\lambda1)=(xy-\lambda1)x$. Intercambiar $x,y$ completa la prueba. La formulación mediante eliminación del cero expresa de forma simétrica exactamente qué parte de los dos espectros coincide.
\end{proof}

\begin{corolario}[Gelfand--Mazur]\label{v3:gelfand-mazur}
Toda álgebra de Banach compleja de división es isomorfa a $\C$ mediante $\lambda\mapsto\lambda1$; con $\norm1=1$, el isomorfismo es isométrico.
\end{corolario}
\begin{proof}
Para $x\in A$ elija $\lambda\in\sigma(x)$. Como todo elemento no nulo es invertible, necesariamente $x-\lambda1=0$. La aplicación indicada es entonces sobreyectiva, además de lineal, multiplicativa e inyectiva.
\end{proof}

\begin{corolario}
Si un álgebra de Banach compleja unitaria $A$ no es isomorfa a $\C$, tiene un divisor topológico de cero no nulo. Si además es conmutativa, tiene un ideal propio no nulo.
\end{corolario}
\begin{proof}
Elija $x$ que no sea múltiplo de $1$. El compacto no vacío $\sigma(x)$ tiene frontera no vacía: un subconjunto de $\C$ con frontera vacía sería simultáneamente abierto y cerrado, y por conexidad sería vacío o todo $\C$, alternativas imposibles. Para $\lambda\in\partial\sigma(x)$ hay $\lambda_n\notin\sigma(x)$ con $\lambda_n\to\lambda$. Así $x-\lambda1$ es un punto no invertible adherente a $G(A)$; el ejercicio~\ref{ej:v3-l21-5} prueba que es divisor topológico de cero. Es no nulo por la elección de $x$.

En el caso conmutativo, si los únicos ideales fueran $0,A$, para cada $x\ne0$ el ideal $xA$ sería $A$ y existiría $y$ con $xy=1$. La conmutatividad haría que $y$ fuera inversa bilateral y Gelfand--Mazur daría $A\simeq\C$, contradicción.
\end{proof}

\begin{observacion}[El ambiente importa]\label{v3:fallo-permanencia}
Si $B\subset A$ es una subálgebra cerrada con la misma identidad y $x\in B$, siempre $\sigma_A(x)\subset\sigma_B(x)$: una inversa en $B$ también lo es en $A$. La igualdad afirmada sin condiciones adicionales en el corolario 9 de la lectura 22 es falsa.

En efecto, tome $A=C(\mathbb T)$ y sea $B$ la clausura uniforme en $\mathbb T$ de los polinomios en $z$ con potencias no negativas. Es una subálgebra cerrada unitaria. Para $u(z)=z$, $\sigma_A(u)=\mathbb T$. En cambio $\sigma_B(u)=\overline{\mathbb D}$. Para justificarlo sin presuponer una representación de $B$, note que el principio del máximo da $|p(w)|\le\norm p_{\mathbb T}$ cuando $|w|\le1$. Por ello la evaluación $p\mapsto p(w)$ se extiende a un homomorfismo $\chi_w:B\to\C$. Si $u-w$ tuviera inversa en $B$, aplicar $\chi_w$ daría $0=1$. Luego todo el disco cerrado pertenece al espectro; la inclusión opuesta se sigue de $\norm u=1$. Más adelante probaremos que la igualdad sí se cumple para subálgebras $C^*$ unitarias.
\end{observacion}

\begin{teorema}[Semicontinuidad superior del espectro]\label{v3:semicontinuidad-espectro}
Si $x_n\to x$ en norma, para cada $\eps>0$ existe $N$ tal que
$\sigma(x_n)\subset\sigma(x)+B(0,\eps)$ cuando $n\ge N$.
\end{teorema}
\begin{proof}
Sea $U=\sigma(x)+B(0,\eps)$. El resolvente $R_x$ está uniformemente acotado en $\C\setminus U$: fuera de un disco grande esto se sigue de la cota geométrica y dentro del disco, de continuidad sobre el compacto que queda al quitar $U$. Sea $M$ una cota positiva. Si $\norm{y-x}<1/M$ y $\lambda\notin U$, entonces
\[
 \lambda1-y=(\lambda1-x)\bigl(1-R_x(\lambda)(y-x)\bigr)
\]
es invertible por Neumann. Por tanto $\sigma(y)\subset U$, y finalmente se aplica la convergencia $x_n\to x$.
\end{proof}

\section{Ejercicios sobre el espectro}

\begin{ejercicio}\label{ej:v3-l22-1}
Si $p$ es un polinomio complejo, pruebe $\sigma(p(x))=p(\sigma(x))$.
\end{ejercicio}


\begin{ejercicio}\label{ej:v3-l22-2}
Para $x\in A$, $y\in B$, pruebe $\sigma_{A\times B}((x,y))=\sigma_A(x)\cup\sigma_B(y)$.
\end{ejercicio}


\begin{ejercicio}\label{ej:v3-l22-3}
Pruebe que $\lambda\mapsto(x-\lambda1)^{-1}$ es diferenciable en el resolvente y que su derivada de Fréchet es continua.
\end{ejercicio}


\begin{ejercicio}\label{ej:v3-l22-4}
En $C^1([0,1])$ con $\norm f=\norm f_\infty+\norm{f'}_\infty$, calcule el radio espectral de $f(t)=t$.
\end{ejercicio}


\begin{ejercicio}\label{ej:v3-l22-5}
Pruebe $r(xy)=r(yx)$.
\end{ejercicio}


\begin{ejercicio}\label{ej:v3-l22-6}
Si $A$ es conmutativa, pruebe
\[
 \sigma(x+y)\subset\sigma(x)+\sigma(y),\qquad
 \sigma(xy)\subset\sigma(x)\sigma(y),
\]
y dé ejemplos de contenciones propias.
\end{ejercicio}


\section{Evaluar funciones holomorfas en un elemento}\label{lec:v3-23}
\fuentelectura{23}{El cálculo funcional holomorfo}

Una función holomorfa cerca del espectro contiene suficiente información para definir $f(x)$. No se requiere que esté definida en todo $\C$, ni que su dominio sea simplemente conexo. Esta precisión permite, por ejemplo, usar $f(z)=1/z$ cuando el espectro evita cero.

\begin{definicion}
Sea $K=\sigma(x)$. Denotamos por $\mathcal O(K)$ el álgebra de \emph{gérmenes holomorfos cerca de $K$}: dos funciones definidas en abiertos que contienen a $K$ representan el mismo germen si coinciden en alguna vecindad de $K$. Para $f$ holomorfa en $U\supset K$, elija una región poligonal acotada $\Omega$ con
$K\subset\Omega\subset\overline\Omega\subset U$ y $\partial\Omega\cap K=\varnothing$, y oriente positivamente su frontera completa $\Gamma$ (las fronteras de los agujeros tienen orientación negativa). Definimos
\[
 f(x)=\frac1{2\pi i}\int_\Gamma f(z)R_x(z)\,dz.
\]
\end{definicion}
\index{cálculo funcional!holomorfo}

Tal región se construye con una cuadrícula suficientemente fina: tome una unión finita de cuadrados que cubra $K$, añada una capa de cuadrados vecinos y conserve su clausura dentro de $U$. La frontera queda separada de $K$. Puede haber varias componentes y agujeros; la orientación completa asegura índice uno sobre $K$ y cero fuera de $\overline\Omega$. La integral existe porque el integrando es continuo en una curva compacta. Dos elecciones dan el mismo valor: se elige una tercera región más pequeña alrededor de $K$, contenida en ambas, y se aplica Cauchy a la diferencia de fronteras, donde el resolvente es holomorfo. Este argumento prueba también que sólo importa el germen de $f$.

\begin{teorema}[Cálculo holomorfo]\label{v3:calculo-holomorfo}
La aplicación $H_x:\mathcal O(\sigma(x))\to A$, $f\mapsto f(x)$, es un homomorfismo unital. Envía $z$ a $x$, es compatible con la composición y satisface:
\begin{enumerate}[label=(\roman*)]
\item $(\alpha f+\beta g)(x)=\alpha f(x)+\beta g(x)$ y $(fg)(x)=f(x)g(x)$;
\item si $g$ es holomorfa cerca de $f(\sigma(x))$, entonces $(g\circ f)(x)=g(f(x))$;
\item si $f_n\to f$ uniformemente en compactos de un mismo abierto que contiene a $\sigma(x)$, entonces $f_n(x)\to f(x)$ en norma.
\end{enumerate}
Además,
\[
 \sigma(f(x))=f(\sigma(x)).
\]
\end{teorema}
\begin{proof}
Para $1$ y $z$, podemos integrar sobre un círculo grande, pues ambas funciones son enteras. La expansión geométrica de $R_x$ y las integrales de los monomios dan $1(x)=1$ y $z(x)=x$. La linealidad es inmediata. Para el producto elija dos fronteras anidadas $\Gamma_1,\Gamma_2$, la primera dentro de la segunda, ambas en un dominio común de $f,g$. La identidad resolvente
\[
 R_x(z)R_x(w)=\frac{R_x(w)-R_x(z)}{z-w}
\]
se verifica multiplicando por $(z1-x)(w1-x)$. Integrando primero en $z$ en el término con $R_x(w)$ se obtiene cero, porque $w$ está fuera de la región interior de $\Gamma_1$. Integrando primero en $w$ en el término con $R_x(z)$, la fórmula de Cauchy da $g(z)$. Fubini para funciones continuas sobre las curvas justifica el cambio de orden. Queda
\[
 f(x)g(x)=\frac1{2\pi i}\int_{\Gamma_1}f(z)g(z)R_x(z)\,dz=(fg)(x).
\]

Probemos ahora el mapeo espectral, antes de usarlo para la composición. Si $\mu\notin f(\sigma(x))$, la función $1/(f-\mu)$ es holomorfa en una vecindad del espectro. El homomorfismo ya probado produce una inversa de $f(x)-\mu1$. Así $\sigma(f(x))\subset f(\sigma(x))$. Recíprocamente, para $\lambda\in\sigma(x)$, defina
\[
 h(z)=\frac{f(z)-f(\lambda)}{z-\lambda},\quad h(\lambda)=f'(\lambda).
\]
La singularidad es removible, y $f(x)-f(\lambda)1=(x-\lambda1)h(x)$. Los dos factores conmutan, pues pertenecen a la imagen de un álgebra conmutativa. Si el producto fuera invertible, también lo sería $x-\lambda1$, contradicción. Se obtiene la inclusión inversa.

Para la composición, tome una frontera $\Delta$ en el dominio de $g$ que encierre $f(\sigma(x))$. Se puede reducir la vecindad de $\sigma(x)$ para que $f$ de su clausura esté dentro de esa región y lejos de $\Delta$. Para $w\in\Delta$, el homomorfismo muestra
\[
 (w1-f(x))^{-1}=\frac1{2\pi i}\int_\Gamma
 \frac{R_x(z)}{w-f(z)}\,dz.
\]
Sustituya esta fórmula en la integral que define $g(f(x))$. Fubini y Cauchy en la variable $w$ convierten la integral interior en $g(f(z))$, y queda exactamente $(g\circ f)(x)$.

Por último, sobre una frontera fija común,
\[
 \norm{f_n(x)-f(x)}\le
 \frac{\operatorname{long}(\Gamma)}{2\pi}
 \max_{z\in\Gamma}\norm{R_x(z)}\,
 \max_{z\in\Gamma}|f_n(z)-f(z)|\longrightarrow0.
\]
La estimación sólo necesita convergencia uniforme en la curva elegida; la convergencia uniforme en compactos la garantiza.
\end{proof}

\begin{proposicion}[Derivada de una función de un elemento]\label{v3:derivada-calculo}
Si $f$ es holomorfa en $U$, la aplicación $x\mapsto f(x)$ es diferenciable en el abierto $\{x:\sigma(x)\subset U\}$. Su derivada es
\[
 Df_A(x)[h]=\frac1{2\pi i}\int_\Gamma f(z)R_x(z)hR_x(z)\,dz.
\]
\end{proposicion}
\begin{proof}
La apertura del dominio se sigue del teorema~\ref{v3:semicontinuidad-espectro}. Fije $\Gamma$ y $M=\max_\Gamma\norm{R_x(z)}$. Para $M\norm h<1/2$, la serie geométrica uniformemente convergente sobre $\Gamma$ da
\[
 R_{x+h}(z)=R_x(z)+R_x(z)hR_x(z)+E_z(h),
 \qquad \norm{E_z(h)}\le2M^3\norm h^2.
\]
Integre multiplicando por $f(z)$. La parte lineal es acotada por
$\operatorname{long}(\Gamma)\max_\Gamma|f|M^2\norm h/(2\pi)$; el resto es $O(\norm h^2)$ uniformemente, y por tanto es $o(\norm h)$. Que la integral lineal tienda a cero no habría bastado por sí solo: es esta cota del resto la que demuestra diferenciabilidad. En general la derivada no es $f'(x)h$, salvo cuando $h$ conmuta con $x$.
\end{proof}

\section{Ejercicios sobre el cálculo holomorfo}

\begin{ejercicio}\label{ej:v3-l23-1}
Si $T\in\Lcal(V)$ y $Tv=\lambda v$ con $v\ne0$, pruebe que $f(T)v=f(\lambda)v$ para toda $f$ holomorfa cerca de $\sigma(T)$.
\end{ejercicio}


\begin{ejercicio}\label{ej:v3-l23-2}
Si $f$ toma el valor cero en algún punto de $\sigma(x)$, pruebe que $f(x)$ no es invertible.
\end{ejercicio}


\begin{ejercicio}\label{ej:v3-l23-3}
Sean $A,B$ álgebras de Banach complejas unitarias y $T:A\to B$ un homomorfismo \emph{continuo y unital}. Pruebe $f(Tx)=T(f(x))$ para $f$ holomorfa cerca de $\sigma_A(x)$.
\end{ejercicio}


\begin{ejercicio}\label{ej:v3-l23-4}
Si $e^x=1$, demuestre $\sigma(x)\subset2\pi i\Z$.
\end{ejercicio}


\begin{ejercicio}\label{ej:v3-l23-5}
Sea $f(z)=\sum_{n\ge0}a_nz^n$ con radio de convergencia $R>r(x)$. Pruebe $f(x)=\sum_{n\ge0}a_nx^n$.
\end{ejercicio}


\begin{ejercicio}\label{ej:v3-l23-6}
Para $T\in\Lcal(V)$, formule y pruebe la compatibilidad del cálculo holomorfo con el adjunto en el dual. Distinga esta operación del adjunto de Hilbert.
\end{ejercicio}


\begin{ejercicio}\label{ej:v3-l23-7}
Sea $I$ un ideal de $A$, no necesariamente cerrado, $x\in I$, y sea $f$ holomorfa en una vecindad de $\sigma(x)\cup\{0\}$, con $f(0)=0$. Pruebe $f(x)\in I$.
\end{ejercicio}


La función resolvente ha convertido una cuestión algebraica ---la existencia de inversas--- en una herramienta analítica. En un álgebra conmutativa veremos que el espectro puede describirse mediante homomorfismos escalares. En las álgebras $C^*$, la involución permitirá extender después el cálculo de funciones holomorfas a todas las funciones continuas del espectro. Para otras presentaciones y desarrollos del cálculo holomorfo pueden consultarse \cite{conway1990,rudin1991,murphy1990}.
